% --- Tabel Deskripsi UC-03 ---
\begin{table}[H]
    \centering
    \caption{Deskripsi Use Case Melakukan \textit{Peer Assessment} (UC-03)}
    \label{tbl:desc_uc03}
    \begin{tabular}{|p{0.25\textwidth}|p{0.70\textwidth}|}
        \hline
        \textbf{Nama Use Case} & \textbf{Melakukan \textit{Peer Assessment}} \\
        \hline
        \textbf{Aktor} & Siswa \\
        \hline
        \textbf{FR Terkait} & FR-03, FR-09 \\
        \hline
        \textbf{Deskripsi} & Siswa mengevaluasi pekerjaan rekan sejawat menggunakan rubrik terstruktur. Aktivitas ini dipicu oleh sistem untuk memecahkan isolasi kompetensi. \\
        \hline
        \textbf{Kondisi sebelum} & Sistem mendeteksi stagnasi pembelajaran; Algoritma \textit{Re-ranking} telah menemukan pasangan heterogen. \\
        \hline
        \textbf{Kondisi sesudah} & Skor Elo \textit{Assessee} diperbarui (berdasarkan nilai tugas); Skor Elo \textit{Assessor} diperbarui (berdasarkan kualitas reviu). \\
        \hline
        \textbf{Main Flow} & \begin{enumerate}[noitemsep,topsep=0pt,leftmargin=*]
            \item Sistem menampilkan jawaban rekan dan rubrik penilaian (Bloom).
            \item Siswa memberikan penilaian kuantitatif per aspek.
            \item Siswa menuliskan komentar kualitatif.
            \item Sistem menghitung kualitas umpan balik (keyword analysis).
            \item Sistem menyimpan hasil penilaian dan memperbarui model kedua siswa.
        \end{enumerate} \\
        \hline
    \end{tabular}
\end{table}